% THETA3-B17 contribution by /root/definite_bridge, 21 November 2026.
\subsection{THETA3-B17: the actual wedge admits an adjoint double lift and last-step seesaw}\label{B17}

The compact-quotient condition in the analytic last-step theorem is
strictly weaker than the totally definite condition in the arithmetic
\texttt{MainX} family. This distinction permits a complete analytic
specialization for the actual multiplicity plane of A7. The second
lift can be its adjoint, with terminal space equal to the original
rank-three space.

\paragraph{Fixed data and the finite-rank operator.}
Use the rank-three case of A7, with its exact scalar automorphic
realization of the canonical coefficient. Thus $V$ has signature
$(2,1)$ at $\mu$ and $(3,0)$ elsewhere, $[F:\Q]\geq2$, and the
multiplicity plane $P=P_1\perp P_2$ is definite at $\mu$.
Write
\[
 G=\U(V),\qquad S=\U(P).
\]
Both adelic quotients $[G]$ and $[S]$ are compact. For $V$, use
the definite component at a place other than $\mu$; for $P$, use
its definite component at $\mu$. A rational isotropic vector would
remain a nonzero isotropic vector in that definite completion.
The quotient compactness is the anisotropic case of the usual
arithmetic compactness theorem already used in A7.

For the nonvanishing assertion about a particular wedge, first
choose its rational negative line $L$ as in B12 and rebase its
canonical coefficient frame at $L$. Make this choice before
fixing the kernel and the operators below. The original
coefficient at an unrelated base point need not be nonzero,
and its maximal compact subgroup need not contain $\U(L^\perp)$.
After this rebasing the following data remain common to all
source characters under consideration.

Fix the additive character, the two splitting characters
$(\chi_V,\chi_P)$, the theta input $\Phi$, finite levels, the
canonical coefficient frame, and positive Haar quotient measures
$dg,dh$. Let
\[
 K(g,h)=\theta_{V,P}(g,h;\Phi).
\]
The coefficient identification here is the one in A7: after
evaluation on the fixed canonical frame, this is an actual scalar
theta kernel, not a separately chosen function having the same
cohomology class. More explicitly, let $\Psi$ be the holomorphic
top-degree relative Lie algebra cochain and choose a unit vector
$\kappa$ in its one-dimensional canonical tangent-frame line.
Its scalar Schwartz coefficient is $\Phi=\Psi(\kappa)$; evaluation
of the theta cochain on the frame translated by $g$ is
$\theta(\omega(g,h)\Phi)$. This fixes the archimedean coefficient
pairing before any adjoint is taken. A fixed conversion between
the Hermitian/skew-Hermitian presentation and the archive's
Gram-matrix presentation is part of the oscillator identification;
the subsequent dagger is applied to this very coefficient.
By A7, $K(\,\cdot\,,h)$ belongs to one
finite-dimensional space $\mathcal E$ of holomorphic canonical
coefficients for every $h$. Its smooth automorphy and the
compactness of the two quotients imply that $K$ is bounded on
$[G]\times[S]$.

Equip $\mathcal E$ with the Petersson inner product inherited from
$L^2([G],dg)$, conjugate linear in the first variable. Define
\begin{equation}
\begin{aligned}
 \mathcal T:L^2([S],dh)&\longrightarrow\mathcal E,\\
 (\mathcal T u)(g)&=\int_{[S]}K(g,h)u(h)\,dh,\\
 \mathcal E_0&=\operatorname{im}\mathcal T.
\end{aligned}
\tag{B17a}
\end{equation}
This is the archive's bilinear \texttt{liftS} convention. Boundedness
follows from Cauchy--Schwarz on the product of the two finite-measure
quotients. Its range is finite dimensional and therefore closed.
Moreover $K(\,\cdot\,,h)\in\mathcal E_0$ for every $h$:
if $f\in\mathcal E$ is orthogonal to the range, then its
continuous coefficient $\langle f,K(\,\cdot\,,h)\rangle_G$
vanishes almost everywhere by testing against all $u$, hence
everywhere. Taking the orthogonal complement proves the claim.
All the A7 character integrals with these fixed data lie in
$\mathcal E_0$: A7 writes the integral for a character $\pi$ as
$f_\pi=\mathcal T(\overline{v_\pi})$.

\paragraph{The adjoint is another genuine theta lift.}
The Hilbert-space adjoint is
\[
 (\mathcal T^*f)(h)
   =\int_{[G]}\overline{K(g,h)}f(g)\,dg.
 \tag{B17b}
\]
Indeed, expanding the inner products and using Fubini gives
$\langle\mathcal T u,f\rangle_G
 =\langle u,\mathcal T^*f\rangle_S$.
In particular $\mathcal T^*f$ is a smooth automorphic function
when $f\in\mathcal E$.

This adjoint kernel is a theta kernel for the exchanged pair with
the same underlying spaces. Since the canonical compact type is
one dimensional, conjugating the chosen scalar coefficient is
exactly the adjoint operation on that coefficient line; there is
no unrecorded sum over compact-type coordinates. Any volume of
the compact fibre is already included in $dg$ and in the
inherited Petersson norm. In the recovered source let
\[
 d_{V,P}(x)=G_V^{-1}x^*G_P,\qquad
 \Phi^\dagger(w)=\overline{\Phi(d_{P,V}(w))},
 \qquad \omega^\dagger=\operatorname{conjAdjoint}(\omega).
\]
Here $G_V,G_P$ are the Gram matrices and $x^*$ is conjugate
transpose. The exact source rule is
\begin{equation}
\begin{gathered}
 \omega\text{ has data }(V,P;\chi_V,\chi_P)
 \quad\Longrightarrow\quad
 \omega^\dagger\text{ has data }(P,V;\chi_P,\chi_V),\\
 \theta_{P,V}^{\dagger}(h,g;\Phi^\dagger)
       =\overline{\theta_{V,P}(g,h;\Phi)}.
\end{gathered}
\tag{B17c}
\end{equation}
There is no inverse on the displayed exchanged splitting
characters in this archive convention.

For precision,
\texttt{Development/\allowbreak{}AbelianVarietyTheory/\allowbreak{}BlockTheory/\allowbreak{}TwistedDiscreteSeries/\allowbreak{}ConjugateAdjointDagger.lean}, lines
112--124, defines the rational involution and the conjugated
exchanged representation. Lines 144--156 give its
antisymplectic, theta-function, and Weil-data clauses.
The following proof files have the full directory prefix
\texttt{Development/\allowbreak{}AbelianVarietyTheory/\allowbreak{}Proofs/\allowbreak{}DecoratedLatticeProofs/}.
The actual proof \texttt{ConjWeilDatumViaAdjointExchangedPair.lean},
theorem \texttt{isWeilDatumO\_conjAdjoint}, proves the exchanged
Weil data, including the split-place normalization.
\texttt{HermitianAdjointSymplecticActionThetaHeisenberg.lean}
proves the rational, antisymplectic, and theta-function
transport. The complete rule is proved by
\texttt{WeilConjugationAdjointRuleStandardChar.lean},
theorem \texttt{conjAdjoint\_holds}, and exposed by
\texttt{GlobalTheory.T.weilConjugationAdjointRule\_standard\_char}
in its statement file.

One can also check the kernel equality directly. The involution
$d$ carries the rational lattice bijectively onto the exchanged
rational lattice, reverses the symplectic form, and exchanges
the two group actions. Complex conjugation reverses the additive
character at the same time, so the standard additive character
is retained on the exchanged pair. The defining formula
$\omega^\dagger(h,g)\Phi^\dagger
 =\overline{\omega(g,h)\Phi}\circ d_{P,V}$
then gives (B17c) at zero, equivalently after summing the
rational Schwartz model. Thus (B17b) is the bilinear theta lift
from $V$ to $P$ with input $f$ and kernel $\Phi^\dagger$.

\paragraph{An exact double lift, common to all source characters.}
On $\mathcal E_0$ put
\[
 Q=\mathcal T\mathcal T^*.
\]
This is positive definite. In fact
$\langle f,Qf\rangle_G=\|\mathcal T^*f\|_S^2$;
if it vanishes, then $f$ is orthogonal to
$\operatorname{im}\mathcal T=\mathcal E_0$, and hence $f=0$.
Consequently, for every $f\in\mathcal E_0$,
\begin{equation}
\begin{gathered}
 F_f=Q^{-1}f,\qquad
 f=\mathcal T\mathcal T^*F_f\\
   =\Theta_{V\leftarrow P}
          \bigl(\Theta_{P\leftarrow V}(F_f;\Phi^\dagger);
                \Phi\bigr).
\end{gathered}
\tag{B17d}
\end{equation}
All functions in this identity are continuous automorphic
functions on compact adelic quotients. It requires neither a new
global nonvanishing theorem nor a rank-one reverse descent:
the second source is $T=V$, of rank three.

The same $Q$ works for every A7 source character with the fixed
spaces, kernel, splittings, levels, coefficient frame, and
measures. Thus (B17d) is a common linear operator identity on the
whole space $\mathcal E_0$, rather than a normalization chosen
separately for each vector or character. Changing the complete
types may change $P$; changing the Schwartz input may change
$Q$. Uniformity across those changes is not asserted.

The exact mixed-pairing identities are
\begin{equation}
\begin{aligned}
 \langle f_1,f_2\rangle_G
    &=\langle\mathcal T^*Q^{-1}f_1,
                        \mathcal T^*f_2\rangle_S\\
    &=\langle\mathcal T^*Q^{-1/2}f_1,
                        \mathcal T^*Q^{-1/2}f_2\rangle_S .
\end{aligned}
\tag{B17e}
\end{equation}
They follow by adjunction and self-adjointness of $Q$.
The positive square root is defined by the finite-dimensional
spectral theorem. In particular, the second identity is an
isometry from $\mathcal E_0$ to a space of actual reverse theta
lifts on the original anisotropic $P$. The conversion from this
Petersson pairing to the geometric canonical cup-product
pairing uses the fixed coefficient frame and quotient measure;
it is the same conversion for every $f_\pi$.

For completeness, when $\mathcal E_0\ne0$, fix an orthonormal basis $f_1,\ldots,f_r$ of
$\mathcal E_0$ and write $K(g,h)=\sum_i f_i(g)e_i(h)$. The matrix
of $Q$ is the positive Gram matrix
\[
 Q_{ij}=\int_{[S]}e_i(h)\overline{e_j(h)}\,dh .
\]
With $t=\operatorname{Tr}Q$ and $\Delta=\det Q>0$, for $r\geq2$
its least eigenvalue satisfies
\[
 \lambda_{\min}(Q)\geq
       \Delta\left(\frac{r-1}{t}\right)^{r-1},
 \qquad
 \|Q^{-1}\|\leq
       \frac{(t/(r-1))^{r-1}}{\Delta}.
 \tag{B17f}
\]
Indeed the product of the other $r-1$ eigenvalues is at most
$(t/(r-1))^{r-1}$ by the arithmetic--geometric mean inequality.
For $r=1$, $\|Q^{-1}\|=\Delta^{-1}$.
Also
$\|\mathcal T^*Q^{-1}f\|_S^2
 =\langle Q^{-1}f,f\rangle_G
 \leq\lambda_{\min}(Q)^{-1}\|f\|_G^2$.
These give an explicit analytic matrix certificate. B7's
arithmetic determinant bound applies only if its arithmetic
normalization hypotheses for this Gram matrix are separately
verified.

\paragraph{Specialization of the actual rank-zero Bessel identity.}
Choose a rational negative line $L\subset V$ at $\mu$ and put
$A=L^\perp$, $H=\U(A)$. Then $A$ is rank two and totally
positive definite. Put
\[
 L_z=-(N_{E/F}\xi_0)^{-1}L,\qquad
 T=V,\qquad X=T\perp L_z .
\]
The identity map between the underlying one-dimensional vector
spaces supplies the exact line relation
$G_L=-N_{E/F}(\xi_0)G_{L_z}$ required by
\texttt{SeeLast}. Use identity rank-zero Witt decompositions
for $V$ and $P$ and the same measures on the identified copies.
The zero-dimensional unipotent measure has mass one.

Choose a splitting $\chi_z$ for $L$ and put
$\chi_A=\chi_V\chi_z^{-1}$. The orthogonal source decomposition
$V=A\perp L$ then has character identity
$\chi_V=\chi_A\chi_z$. The exchanged second pair in (B17c)
has characters $(\chi'_P,\chi_T)=(\chi_P,\chi_V)$, so the
upper pair $(P,X)$ has characters
\[
 (\chi'_P,\chi_T\chi_z^{-1})=(\chi_P,\chi_A),
 \qquad (P,L_z)\text{ has }(\chi_P,\chi_z^{-1}).
\]
These are exactly the character equations in
\texttt{Development/\allowbreak{}AbelianVarietyTheory/\allowbreak{}Definitions/\allowbreak{}LatticeFrames/\allowbreak{}IsQuasiFam.lean},
lines 114--117. As in B15f,
such global splitting characters exist; their selection is
part of the fixed input, with no assertion about prescribed
unramified finite values.

Let $e$ be a continuous automorphic test function on $[H]$, put
$\varphi=\overline e$, and use the rank-zero Bessel period
\[
 B_e(f)=\int_{[H]}f(\iota_A(a))\overline{e(a)}\,da .
\]
In the notation of (B15l), let
$C=C_{\Phi,\Phi^\dagger}=\texttt{lastBi}(\Phi,\Phi^\dagger)$
be the exact contracted two-variable theta input. Its formula
on $X=V\perp L_z$ is
\[
 C(u;v_T,v_z)
     =\Phi^\dagger(v_T)\,
       \Phi\bigl((o.\operatorname{res}_S)^{-1}
                       (u,\operatorname{ofAdj}_{L,P}(v_z))\bigr).
\]
The line map here is the identity chosen above; the source
definition incorporates its Gram matrices and $\xi_0$ in
$\operatorname{ofAdj}$. Define $\mathcal J_e(x)$ to be the
archive's upper double lift
$\operatorname{biLift}_{A\to P\to X}(\overline e;C)(x)$,
with the same $da,dh$ as above. Then the exact identity is
\begin{equation}
 B_e(f)=\int_{[G]}(Q^{-1}f)(t)\,
                 \mathcal J_e(\iota_T(t,1))\,dt
       \qquad(f\in\mathcal E_0).
\tag{B17g}
\end{equation}

\textbf{Proof.}
Set
$\beta_e(h)=\int_{[H]}\overline{e(a)}K(\iota_A(a),h)\,da$.
Substitute (B17d) in the left side and use Fubini on
$[H]\times[S]\times[G]$. The result is
\[
 \int_{[G]}F_f(t)\int_{[S]}
       \overline{K(t,h)}\,\beta_e(h)\,dh\,dt.
\]
All factors are continuous on the respective compact quotients,
so their product is absolutely integrable. Formula (B17c)
identifies the first kernel in the inner integral with
$\theta^\dagger_{P,V}(h,t;\Phi^\dagger)$.
The orthogonal tensor identity and the line relation now give
\[
 \mathcal J_e(\iota_T(t,1))
       =\int_{[S]}\theta^\dagger_{P,V}(h,t;\Phi^\dagger)
                         \beta_e(h)\,dh .
\]
This is precisely the integrand equality of
\texttt{SeeLast}, after integration over $[H]$ and $[S]$.
It proves (B17g).

These are the actual hypotheses used by the recovered analytic
theorems. Both definition files in this paragraph have directory
prefix
\texttt{Development/\allowbreak{}AbelianVarietyTheory/\allowbreak{}Definitions/\allowbreak{}LatticeFrames/}.
\texttt{IsThetaFam.lean}, line 17, defines
\texttt{CoCpt} by a compact set of adelic representatives;
its structures \texttt{BesLev}, \texttt{JacLev}, and
\texttt{Last} ask for this condition, not for definiteness at
every real place. \texttt{IsQuasiFam.lean}, lines 71--75,
states \texttt{LastLaw}, and lines 110--123 state
\texttt{SeeLast}. The complete proofs, in
\texttt{Development/\allowbreak{}AbelianVarietyTheory/\allowbreak{}Proofs/\allowbreak{}DecoratedLatticeProofs/}, are
\texttt{DoubleLiftPeriodRankZeroFubiniCompactQuotient.lean},
lines 136--159, and \texttt{LastSeeSawThetaIdentity.lean},
lines 22--54. The proof above instantiates their analytic
content with $S=P$ and $T=V$. It does not require the
additional \texttt{Comm.Is} hypotheses used to package
\texttt{GlobalTheory.LastRed.main}. \hfill$\square$

\paragraph{What is now obtained, and what remains arithmetic.}
For a nonzero actual A7 wedge, make the B12 choice of rational
$H$ and rebase the coefficient frame before fixing $K,Q$, as
specified above. B12 then supplies an $e$ for which $B_e(f)\ne0$.
Equation (B17g) is therefore a
nonzero actual last-step Bessel seesaw involving that wedge.
It is uniform in the source characters for the fixed data.
The identity uses a common positive operator $Q^{-1}$, not
an assertion that this operator is a scalar or is algebraic.

Equation (B17g) is an identity for one Bessel coefficient.
The entire geometric mixed period is recovered from the
Petersson pairing in (B17e), with the fixed geometric measure
conversion. If a description solely through definite Bessel
coefficients is wanted, apply the finite Bessel frame of A5 to
$\mathcal E_0$, and then apply (B17g) to each of its coefficients.
For an adelic translate in that frame, translate the first
kernel by the same group element; \texttt{LastLaw} permits the
two kernels to be different. This combines actual last-step
seesaws to recover the whole pairing. A single restriction
period is not being identified with the whole ball integral.

The upper space is completely determined here. Since
$N_{E/F}\xi_0$ is a norm, $L\perp L_z$ is a hyperbolic plane:
if $e$ and $e_z$ have Gram entries $a$ and
$-a/N_{E/F}\xi_0$, the vector $e+\xi_0e_z$ is isotropic.
Thus
\[
 X\simeq A\perp\mathbb H,\qquad
 \operatorname{rank}X=4,\qquad
 \operatorname{sig}_\tau X=(3,1)
        \text{ for every real }\tau .
\]
Its Witt index is one, because $A$ is definite at every real
place. The quotient $[X]$ is not being integrated over in
(B17g); evaluation on the embedded compact $[G]$ is sufficient.
This feature causes no convergence gap in the displayed
identity.

The arithmetic \texttt{MainX} family, by contrast, asks for
totally definite spaces, specified arithmetic functionals,
and the common kernel/covolume and terminal mass equations
(B15m). Neither the mixed $P$ nor this rank-four $X$ is replaced
by a definite inner form in the argument. The result closes
the analytic double-lift and last-step comparison for the
actual A7 kernel while keeping those arithmetic matching
conditions separate. In particular it proves no algebraicity
of $Q^{-1}$ and no new projected algebraic cycle.
